\GXNSFBodyText

\noindent\textbf{预期科学结果}\par

\noindent\textbf{理论与方法结果}\par
形成跨媒介职业状态的统一表示方法，刻画平台观测偏差、状态迁移与外部扰动之间的时序关联；给出职业韧性的分量化定义、估计程序和失效边界；形成从公开单案例到同构人物验证、再到多语言方法组件测试的分层研究范式。

\noindent\textbf{预期成果载体}\par

\noindent\textbf{论文、数据与软件}\par
计划形成相关学术论文、事件编码手册和研究报告；发布不含版权原文和个人敏感信息的来源索引、合成示例数据及复现实验脚本；实现用于演示状态分段与韧性敏感性分析的原型工具。正式申请中的论文数量、目标期刊和软件著作权等指标，应由申请人根据真实基础和任务安排据实确定。

\noindent\textbf{预期应用结果}\par

\noindent\textbf{区域数字文化传播分析工具}\par
在广西—东盟公开数字文化传播场景中，考察多语言实体对齐、平台校准和扰动窗口组件的适用性，为区域文化内容的长期传播监测提供方法参考。结果以聚合统计方式呈现，服务于区域文化内容的整体分析。
